\subsection{Primary Decomposition and Spectral Splitting}

\begin{theorem} \label{thm:complex-primary-decomposition-generalized-eigenspaces} \hlt{Primary Decomposition over $\mathbb{C}$}\\
Let $\mathcal{V}\neq\{\mathbf{0}\}$ be finite-dimensional over $\mathbb{C}$,  $T\in\LL(\mathcal{V})$, $m\in\mathbb{Z}_{>0}$. Suppose $\lambda_1,\ldots,\lambda_m$ are the distinct eigenvalues of $T$. Then:
\begin{enumerate}[label=(\roman*)]
\setlength{\itemsep}{0pt}
\item $\mathcal{V}=\mathcal{G}_{\lambda_1}(T)\oplus\cdots\oplus \mathcal{G}_{\lambda_m}(T)$;
\item $(\forall j\in\{1,\ldots,m\})\,(T-\lambda_jI_{\mathcal{V}})|_{\mathcal{G}_{\lambda_j}(T)}$ is nilpotent.
\end{enumerate}
\begin{proof}
Set $n\defeq\dim(\mathcal{V})$. For $j\in\{1,\ldots,m\}$, set $S_j\defeq T-\lambda_jI_{\mathcal{V}}$ and $\mathcal{G}_j\defeq\mathcal{G}_{\lambda_j}(T)$.
\begin{enumerate}[label=(\roman*)]
\setlength{\itemsep}{0pt}
\item[$(i)$] We prove $\mathcal{V}=\mathcal{G}_1\oplus\cdots\oplus\mathcal{G}_m$ by induction on $n$.
\begin{enumerate}[label=(\alph*)]
\setlength{\itemsep}{0pt}
\item \hlt{Base Case}: Suppose $n=1$. Choose $\mathbf{v}\in\mathcal{V}\setminus\{\mathbf{0}\}$ by $\mathcal{V}\neq\{\mathbf{0}\}$. Then $(\mathbf{v})$ is a basis by $n=1$.\\
Choose $\lambda\in\mathbb{C}$ with $T\mathbf{v}=\lambda\mathbf{v}$ by spanning. Then $T=\lambda I_{\mathcal{V}}$ by linearity, so $m=1$, $\lambda=\lambda_1$, and $\mathcal{V}=\mathcal{G}_1$.\\
Assume henceforth that $n>1$ and that (i) holds in every smaller positive dimension.
\item \hlt{Stable Split}: Set $\mathcal{U}\defeq\operatorname{range}(S_1^n)$. By Theorem \ref{thm:kernel-range-stable-direct-sum} and Theorem \ref{thm:generalized-eigenspace-stabilized-kernel}(i), $\mathcal{V}=\mathcal{G}_1\oplus\mathcal{U}$.\\
The subspace $\mathcal{U}$ is $T$-invariant by Proposition \ref{prop:kernel-range-polynomial-in-operator-invariant}(ii).\\
Since $\lambda_1$ is an eigenvalue, $\mathcal{G}_1\neq\{\mathbf{0}\}$ by Definition \ref{def:generalized-eigenvector-and-eigenspace}(i); hence $\dim(\mathcal{U})<n$ by Theorem \ref{thm:direct-sum-iff-dimensions-add}(i).
\item \hlt{Zero Range}: Suppose $\mathcal{U}=\{\mathbf{0}\}$. Then $\mathcal{V}=\mathcal{G}_1$ by item (b). Fix $\lambda\in\operatorname{spec}(T)$ and $\mathbf{v}\in E_{\lambda}(T)\setminus\{\mathbf{0}\}$.\\
Then $0=S_1^n\mathbf{v}=(\lambda-\lambda_1)^n\mathbf{v}$, so $\lambda=\lambda_1$. Thus $\operatorname{spec}(T)=\{\lambda_1\}$, $m=1$, and (i) holds.
\item \hlt{Restricted Spectrum}: Suppose $\mathcal{U}\neq\{\mathbf{0}\}$. Then $T|_{\mathcal{U}}\in\LL(\mathcal{U})$ by item (b), and $\operatorname{spec}(T|_{\mathcal{U}})\neq\emptyset$ by Theorem \ref{thm:complex-operator-has-eigenvalue}.\\
$\operatorname{spec}(T|_{\mathcal{U}})\subseteq\operatorname{spec}(T)\setminus\{\lambda_1\}$ by restriction and $\mathcal{U}\cap\ker(S_1)\subseteq\mathcal{U}\cap\mathcal{G}_1=\{\mathbf{0}\}$. Hence $m\geq2$.\\
Fix $j\in\{2,\ldots,m\}$ and choose an eigenvector $\mathbf{v}_j$ for $\lambda_j$ by $\lambda_j\in\operatorname{spec}(T)$.\\
Then $\mathbf{v}_j=(\lambda_j-\lambda_1)^{-n}S_1^n\mathbf{v}_j\in\mathcal{U}$, so $\operatorname{spec}(T|_{\mathcal{U}})=\{\lambda_2,\ldots,\lambda_m\}$.
\item \hlt{Inductive Decomposition}: By items (b) and (d) and the induction hypothesis, $\mathcal{U}=\mathcal{G}_{\lambda_2}(T|_{\mathcal{U}})\oplus\cdots\oplus\mathcal{G}_{\lambda_m}(T|_{\mathcal{U}})$.
\item \hlt{Eigenspace Identification}: Fix $j\in\{2,\ldots,m\}$. By equality of restriction powers, $\mathcal{G}_{\lambda_j}(T|_{\mathcal{U}})\subseteq\mathcal{G}_j$.\\
Conversely, let $\mathbf{v}\in\mathcal{G}_j$ and write $\mathbf{v}=\mathbf{g}+\mathbf{u}$ with $\mathbf{g}\in\mathcal{G}_1$ and $\mathbf{u}\in\mathcal{U}$ by item (b).\\
Choose $r\in\mathbb{Z}_{>0}$ with $S_j^r\mathbf{v}=\mathbf{0}$ by Definition \ref{def:generalized-eigenvector-and-eigenspace}(ii). Then $S_j^r\mathbf{g}\in\mathcal{G}_1$ and $S_j^r\mathbf{u}\in\mathcal{U}$ by Proposition \ref{prop:generalized-eigenspaces-invariant} and item (b).\\
Both terms vanish by $S_j^r\mathbf{g}+S_j^r\mathbf{u}=\mathbf{0}$ and the directness in item (b). Thus $\mathbf{g}\in\mathcal{G}_1\cap\mathcal{G}_j=\{\mathbf{0}\}$ by Theorem \ref{thm:distinct-generalized-eigenvectors-independent}.\\
Hence $\mathbf{v}=\mathbf{u}\in\mathcal{G}_{\lambda_j}(T|_{\mathcal{U}})$. Therefore the generalized eigenspaces are equal, and items (b) and (e) prove (i).
\end{enumerate}
\item[$(ii)$] Fix $j\in\{1,\ldots,m\}$ and set $N_j\defeq S_j|_{\mathcal{G}_j}$. By Theorem \ref{thm:generalized-eigenspace-stabilized-kernel}(ii) and Proposition \ref{prop:generalized-eigenspaces-invariant}, $N_j\in\LL(\mathcal{G}_j)$.\\
By Theorem \ref{thm:generalized-eigenspace-stabilized-kernel}(i), $N_j^n=0$, so $N_j$ is nilpotent by Definition \ref{def:nilpotent-operator}.
\end{enumerate}
\hlt{Structural Commentary}: Stable kernel-range splitting isolates one primary summand; induction decomposes the invariant range.\\
\hlt{Key Insight}: Stable kernel-range splitting isolates one eigenvalue at a time.
\end{proof}
\end{theorem}

Figure \ref{fig:primary-decomposition-generalized-eigenspaces} summarizes the primary summands and the nilpotent part of $T$ on each one.

\begin{illustrationfigure}[H]
  \centering
  \begin{tikzpicture}[
  >=Stealth,
  every node/.style={font=\small},
  block/.style={draw=blue!55!black, fill=blue!6, rounded corners=3pt, minimum width=2.1cm, minimum height=1.05cm, align=center},
  arrow/.style={->, line width=0.85pt, blue!70!black},
  nil/.style={->, line width=0.75pt, red!65!black}
]
  \node[block] (G1) at (0,0) {$G(\lambda_1,T)$\\$T=\lambda_1 I+N_1$};
  \node[block] (G2) at (2.95,0) {$G(\lambda_2,T)$\\$T=\lambda_2 I+N_2$};
  \node[block] (Gm) at (6.16,0) {$G(\lambda_m,T)$\\$T=\lambda_m I+N_m$};

  \node at (4.58,0) {$\cdots$};
  \node[above=0.62cm of G2, font=\small] (V) {$\mathcal{V}=G(\lambda_1,T)\oplus G(\lambda_2,T)\oplus\cdots\oplus G(\lambda_m,T)$};

  \draw[arrow] (V.south west) to[out=-110,in=90] (G1.north);
  \draw[arrow] (V.south) -- (G2.north);
  \draw[arrow] (V.south east) to[out=-70,in=90] (Gm.north);

  \draw[nil] ($(G1.south west)+(0.25,0.18)$) .. controls +(0,-0.55) and +(0,-0.55) .. ($(G1.south east)+(-0.25,0.18)$)
    node[midway, below=5pt] {$N_1^{r_1}=0$};
  \draw[nil] ($(G2.south west)+(0.25,0.18)$) .. controls +(0,-0.55) and +(0,-0.55) .. ($(G2.south east)+(-0.25,0.18)$)
    node[midway, below=5pt] {$N_2^{r_2}=0$};
  \draw[nil] ($(Gm.south west)+(0.25,0.18)$) .. controls +(0,-0.55) and +(0,-0.55) .. ($(Gm.south east)+(-0.25,0.18)$)
    node[midway, below=5pt] {$N_m^{r_m}=0$};
\end{tikzpicture}

  \caption{Primary decomposition: on $\mathcal{G}_{\lambda_j}(T)$, $T=\lambda_j I$ plus nilpotent.}
  \label{fig:primary-decomposition-generalized-eigenspaces}
\end{illustrationfigure}

\begin{corollary} \label{cor:complex-basis-of-generalized-eigenvectors} \hlt{Basis of Generalized Eigenvectors}\\
Let $\mathcal{V}\neq\{\mathbf{0}\}$ be finite-dimensional over $\mathbb{C}$ and $T\in\LL(\mathcal{V})$. Then $\mathcal{V}$ has a generalized-eigenvector basis for $T$.
\begin{proof}
Enumerate the distinct eigenvalues as $\lambda_1,\ldots,\lambda_m$ by Theorem \ref{thm:complex-operator-has-eigenvalue} and Proposition \ref{prop:operator-characteristic-polynomial-degree-roots}(i)--(ii).\\
By Theorem \ref{thm:complex-primary-decomposition-generalized-eigenspaces}(i), $\mathcal{V}=\mathcal{G}_{\lambda_1}(T)\oplus\cdots\oplus\mathcal{G}_{\lambda_m}(T)$.\\
For $j\in\{1,\ldots,m\}$, the space $\mathcal{G}_{\lambda_j}(T)$ has an ordered basis $\beta_j$ by Theorem \ref{thm:generalized-eigenspace-stabilized-kernel}(ii) and Corollary \ref{cor:finite-dimensional-spaces-and-subspaces-have-bases}(i)--(ii).\\
The concatenation $\beta\defeq(\beta_1,\ldots,\beta_m)$ is an ordered basis of $\mathcal{V}$ by Theorem \ref{thm:direct-sum-iff-dimensions-add}(ii).\\
Every vector in $\beta_j$ is nonzero and belongs to $\mathcal{G}_{\lambda_j}(T)$, so it is a generalized eigenvector corresponding to $\lambda_j$ by Definition \ref{def:generalized-eigenvector-and-eigenspace}(i).\\
\hlt{Structural Commentary}: Primary decomposition supplies finite-dimensional summands whose ordered bases concatenate by directness.\\
\hlt{Key Insight}: Every primary-summand basis vector is generalized for its eigenvalue.
\end{proof}
\end{corollary}

\begin{theorem} \label{thm:complex-block-diagonal-generalized-eigenbasis} \hlt{Block Diagonal Form with Upper-Triangular Primary Blocks}\\
Let $\mathcal{V}\neq\{\mathbf{0}\}$ be finite-dimensional over $\mathbb{C}$ and let $T\in\LL(\mathcal{V})$.\\
Let $m\in\mathbb{Z}_{>0}$, let $(\lambda_j)_{j=1}^{m}$ list its distinct eigenvalues, and set $d_j\defeq\dim\mathcal{G}_{\lambda_j}(T)$ for $j\in\{1,\ldots,m\}$.\\
Then there are an ordered basis $\beta$ of $\mathcal{V}$ and matrices $A_j\in\operatorname{Mat}(d_j,d_j,\mathbb{C})$, $j\in\{1,\ldots,m\}$, such that $[T]_{\beta\leftarrow\beta}=A_1\oplus\cdots\oplus A_m$.\\
Each $A_j$ is upper triangular with diagonal $(\lambda_j,\ldots,\lambda_j)$.
\begin{proof}
\begin{enumerate}[label=(\roman*)]
\setlength{\itemsep}{0pt}
\item \hlt{Primary Summands}: Set $\mathcal{U}_j\defeq\mathcal{G}_{\lambda_j}(T)$ and $N_j\defeq(T-\lambda_jI_{\mathcal{V}})|_{\mathcal{U}_j}$ for $j\in\{1,\ldots,m\}$.\\
By Theorem \ref{thm:complex-primary-decomposition-generalized-eigenspaces}(i)--(ii), $\mathcal{V}=\mathcal{U}_1\oplus\cdots\oplus\mathcal{U}_m$ and each $N_j\in\LL(\mathcal{U}_j)$ is nilpotent.
\item \hlt{Nilpotent Bases}: Fix $j\in\{1,\ldots,m\}$. Choose ordered basis $\beta_j$; $[N_j]_{\beta_j}$ strictly upper triangular by item (i), Theorem \ref{thm:nilpotent-strictly-upper-triangular-matrix}.\\
Set $A_j\defeq[T|_{\mathcal{U}_j}]_{\beta_j\leftarrow\beta_j}$. By Proposition \ref{prop:matrix-representation-respects-linear-combinations}, $A_j=\lambda_jI_{d_j}+[N_j]_{\beta_j}$, so $A_j$ is upper triangular with diagonal $(\lambda_j,\ldots,\lambda_j)$.
\item \hlt{Adapted Basis}: The concatenation $\beta\defeq(\beta_1,\ldots,\beta_m)$ is an ordered basis of $\mathcal{V}$ by Theorem \ref{thm:direct-sum-iff-dimensions-add}(ii).
\item \hlt{Block Verification}: By items (i)--(iii) and Definition \ref{def:matrix-representation-of-linear-map}, $[T]_{\beta\leftarrow\beta}=A_1\oplus\cdots\oplus A_m$.
\end{enumerate}
\hlt{Structural Commentary}: Primary summands determine the blocks; nilpotent-adapted bases determine their triangular shape.\\
\hlt{Key Insight}: On each generalized eigenspace, $T$ is eigenvalue plus nilpotent.
\end{proof}
\end{theorem}

Figure \ref{fig:primary-block-diagonal-form} illustrates the block form obtained from bases adapted to the primary summands.

\begin{illustrationfigure}[H]
  \centering
  \begin{tikzpicture}[
  every node/.style={font=\small},
  mat/.style={draw=gray!60, fill=gray!5, rounded corners=2pt, inner xsep=5pt, inner ysep=4.5pt, align=center},
  block/.style={draw=blue!60!black, fill=blue!7, minimum width=1.48cm, minimum height=0.53cm, align=center},
  arrow/.style={->, line width=0.85pt, blue!70!black}
]
  \node[mat] (M) at (0,0) {
    $\begin{bmatrix}
      A_1 & & 0\\
      & \ddots & \\
      0 & & A_m
    \end{bmatrix}$};

  \node[block] (A1) at (3.26,0.72) {$A_1$};
  \node[block] (Aj) at (3.26,0) {$A_j$};
  \node[block] (Am) at (3.26,-0.72) {$A_m$};

  \draw[arrow] (M.east) -- (A1.west);
  \draw[arrow] (M.east) -- (Aj.west);
  \draw[arrow] (M.east) -- (Am.west);
\end{tikzpicture}

  \caption{Primary blocks: $A_j=\lambda_jI+N_j$, with $N_j$ strictly upper triangular.}
  \label{fig:primary-block-diagonal-form}
\end{illustrationfigure}

\begin{proposition} \label{prop:eigenvalue-algebraic-multiplicity-generalized-eigenspace} \hlt{Algebraic Multiplicity from the Generalized Eigenspace}\\
Let $\mathcal{V}\neq\{\mathbf{0}\}$ be finite-dimensional over $\mathbb{C}$ and let $T\in\LL(\mathcal{V})$.\\
Then $(\forall\lambda\in\operatorname{spec}(T))$, the multiplicity of $\lambda$ as a root of $q_T$ is $\dim\mathcal{G}_{\lambda}(T)=\dim\ker((T-\lambda I_{\mathcal{V}})^{\dim(\mathcal{V})})$.
\begin{proof}
Fix $\lambda\in\operatorname{spec}(T)$. By Theorem \ref{thm:complex-block-diagonal-generalized-eigenbasis}, $\exists$ ordered basis $\beta$,  $A\defeq[T]_{\beta\leftarrow\beta}$ upper triangular, diagonal has $\lambda$ exactly $\dim\mathcal{G}_{\lambda}(T)$ times.\\
By Definition \ref{def:operator-characteristic-polynomial} and Proposition \ref{prop:upper-triangular-characteristic-polynomial}(iii), the multiplicity of $\lambda$ as a root of $q_T$ is $\dim\mathcal{G}_{\lambda}(T)$.\\
The identity $\mathcal{G}_{\lambda}(T)=\ker((T-\lambda I_{\mathcal{V}})^{\dim(\mathcal{V})})$ holds by Theorem \ref{thm:generalized-eigenspace-stabilized-kernel}(i).\\
\hlt{Structural Commentary}: Primary-block diagonal entries convert root multiplicity into generalized-eigenspace dimension.\\
\hlt{Key Insight}: Algebraic multiplicity is the dimension of the generalized eigenspace.
\end{proof}
\end{proposition}

Figure \ref{fig:operator-characteristic-polynomial-primary-blocks} records how the primary-summand dimensions determine the characteristic-polynomial exponents.

\begin{illustrationfigure}[H]
  \centering
  \begin{tikzpicture}[
  >=Stealth,
  every node/.style={font=\small},
  block/.style={draw=blue!60!black, fill=blue!6, rounded corners=2pt, align=center, inner xsep=4.5pt, inner ysep=4pt, minimum width=1.7cm},
  factor/.style={draw=dartmouthgreen!70!black, fill=dartmouthgreen!8, rounded corners=2pt, align=center, inner xsep=4.5pt, inner ysep=4pt, minimum width=1.7cm},
  result/.style={draw=gray!60, fill=gray!6, rounded corners=2pt, align=center, inner xsep=5pt, inner ysep=4.5pt},
  arr/.style={->, line width=0.7pt, draw=black!70}
]
  \node[block] (G1) at (-2.55,0) {$G(\lambda_1,T)$};
  \node (Gdots) at (0,0) {$\cdots$};
  \node[block] (Gm) at (2.55,0) {$G(\lambda_m,T)$};

  \node[factor] (F1) at (-2.55,-1.08) {$(z-\lambda_1)^{d_1}$};
  \node (Fdots) at (0,-1.08) {$\cdots$};
  \node[factor] (Fm) at (2.55,-1.08) {$(z-\lambda_m)^{d_m}$};

  \node[result] (q) at (0,-2.10) {$q_T(z)=\prod_{j=1}^{m}(z-\lambda_j)^{d_j}$};

  \draw[arr] (G1) -- (F1);
  \draw[arr] (Gm) -- (Fm);
  \draw[arr] (F1.south east) -- (q.north west);
  \draw[arr] (Fm.south west) -- (q.north east);
\end{tikzpicture}

  \caption{Each primary summand contributes $(z-\lambda_j)^{d_j}$, where $d_j=\dim \mathcal{G}_{\lambda_j}(T)$.}
  \label{fig:operator-characteristic-polynomial-primary-blocks}
\end{illustrationfigure}

\begin{proposition} \label{prop:sum-of-complex-eigenvalue-multiplicities} \hlt{Sum of Multiplicities}\\
Let $\mathcal{V}\neq\{\mathbf{0}\}$ be a finite-dimensional complex vector space and let $T\in\LL(\mathcal{V})$.\\
For $\lambda\in\operatorname{spec}(T)$, let $a(\lambda)$ be its multiplicity as a root of $q_T$. Then $\sum\limits_{\lambda\in\operatorname{spec}(T)}a(\lambda)=\dim(\mathcal{V})$.
\begin{proof}
By Theorems \ref{thm:complex-primary-decomposition-generalized-eigenspaces}(i) and \ref{thm:direct-sum-iff-dimensions-add}(i), $\dim(\mathcal{V})=\sum\limits_{\lambda\in\operatorname{spec}(T)}\dim \mathcal{G}_{\lambda}(T)$.\\
By Proposition \ref{prop:eigenvalue-algebraic-multiplicity-generalized-eigenspace}, $\dim\mathcal{G}_{\lambda}(T)=a(\lambda)$ for every $\lambda\in\operatorname{spec}(T)$, so the assertion follows.\\
\hlt{Structural Commentary}: Primary decomposition identifies algebraic multiplicities with summand dimensions.\\
\hlt{Key Insight}: Direct-sum dimensions add to $\dim(\mathcal{V})$.
\end{proof}
\end{proposition}
